% Test-case specification TC-0015 (CTest test TC-0015-members).
\documentclass[letterpaper]{article}

\usepackage[letterpaper,margin=1in]{geometry}
\emergencystretch=3em

% Compile from the repository root so this path resolves to
% the pinned WSP submodule.
\input{wsp/testing/test-case-library.tex}

\begin{document}

\def\TCID{TC-0015}
\def\TCTitle{%
  Member assignment, event sinks and the record-array move}
\def\TCPurpose{%
  Verify \texttt{PutMemVar}, \texttt{PutMemObj}, \texttt{SetMemVar}, the
  Event and NewObj forms, the event-sink reference counts and the
  record-array move.}
\def\TCPriority{\TCPriorityHigh}

\def\TCRequirementRef{%
  REQ-0002, REQ-0003, REQ-0005}
\def\TCDesignRef{%
  Not applicable.}
\def\TCFeatureRef{%
  \texttt{tests/\allowbreak{}vbrun\_\allowbreak{}members\_\allowbreak{}tests.\allowbreak{}cpp},
  \texttt{src/\allowbreak{}helpers\_\allowbreak{}members.\allowbreak{}cpp}}
\def\TCTestTechnique{%
  Requirements-based testing; reference counts observed.}

\def\TCPreconditions{%
  The library and the test program are built
  (\texttt{scripts/\allowbreak{}Invoke-Build.\allowbreak{}ps1}) and the test
  is registered with CTest.}
\def\TCEnvironment{%
  Windows with the toolchain of
  \texttt{scripts/\allowbreak{}Invoke-Build.\allowbreak{}ps1}; the local
  gate builds and runs x64 Debug, CI builds and runs the other
  architectures. The library is built from this repository in
  \texttt{out/\allowbreak{}build}.}
\def\TCAssumptions{%
  Not applicable.}

\def\TCInputData{%
  Variants, objects, event sinks made by the test and arrays of records.}
\def\TCInitialState{%
  No state from earlier runs is needed; the test makes what it needs.}

\def\TCProcedure{%
  \begin{enumerate}
    \item CTest runs \texttt{vbrun\_\allowbreak{}members\_\allowbreak{}tests}.
    \item Assign members of the Variant and object kinds and check the values
      and reference counts.
    \item Add and remove event sinks and check the counts.
    \item Move a record array and check the contents.
    \item CTest records the exit code and, on failure, the program's output
      (\texttt{--output-on-failure}).
  \end{enumerate}}
\def\TCExpectedResult{%
  All checks pass; the program prints \texttt{vbrun members: all checks
  passed} and exits 0.}
\def\TCPostConditions{%
  No file or registry state remains that later tests depend on.}

\TCTemplate

\end{document}
